\uuid{4RlG}
\titre{Limite en l'infini}
\auteur{bodin ; hanani ; mzari}
\organisation{exo7}
\licence{CC BY-NC-SA 4.0}
\datecreate{2020-01-20}
\niveau{L1}
\module{Analyse}
\chapitre{Continuité, limite et étude de fonctions réelles}
\sousChapitre{Limite de fonctions}
\difficulte{3}
\theme{Limites des fonctions réelles, limite en l'infini}
\exo7qcmid{lille-1.limites.022}

\contenu{
\begin{question}
\qclassification{123.03}
\qtags{motcle=limite en l'infini}
Soit $f(x)= xE(\frac{1}{x})$, où $E$ désigne la partie entière. Quelles sont les assertions vraies ?
\begin{answers}

    \good{$\lim_{x\to +\infty} f(x)=0$}

    \bad{$\lim_{x\to +\infty} f(x)=+\infty$}

    \bad{$\lim_{x\to -\infty} f(x)=0$}

    \good{$\lim_{x\to -\infty} f(x)=+\infty$}

\end{answers}
\begin{explanations}
Pour $x>1$,  $E(\frac{1}{x})=0$, donc $f(x)=0$ et donc  $\lim_{x\to +\infty} f(x)=0$.
Pour $x<-1$,  $E(\frac{1}{x})=-1$, donc $f(x)=-x$ et donc  $\lim_{x\to -\infty} f(x)=+\infty$.
\end{explanations}
\end{question}
}
